%%%%%%%%%%%%%%%%%%%%%%%%%%%%%%%%%%%%%%%%%%%%%%%%%%%%%%%%%%%%
% 2.3 ALGEBRAIC NUMBER THEORY
% The Composition of Advanced Mathematics
%%%%%%%%%%%%%%%%%%%%%%%%%%%%%%%%%%%%%%%%%%%%%%%%%%%%%%%%%%%%

\section{Algebraic Number Theory}
\label{sec:algebraic-number-theory}

\prerequisite{
Elementary Number Theory from \cref{sec:elementary-number-theory}, together
with Abstract Algebra, Ring Theory, Field Theory, and basic familiarity with
polynomials. Some later topics connect naturally with Galois Theory,
Commutative Algebra, and Analysis.
}

\learningobjectives{
By the end of this section, the reader should be able to:
\begin{itemize}
    \item explain the purpose of algebraic number theory;
    \item distinguish rational, irrational, algebraic, and transcendental numbers;
    \item determine whether a given complex number is algebraic;
    \item work with minimal polynomials and algebraic degree;
    \item understand number fields as finite extensions of \(\Q\);
    \item identify algebraic integers;
    \item explain why the ordinary integers must be enlarged inside number fields;
    \item work with quadratic number fields;
    \item describe rings of integers in basic quadratic fields;
    \item compute traces and norms;
    \item distinguish element factorization from ideal factorization;
    \item explain failure of unique factorization in rings such as \(\Z[\sqrt{-5}]\);
    \item define ideals and prime ideals in number rings;
    \item understand fractional ideals at an introductory level;
    \item state the unique factorization theorem for ideals in rings of integers;
    \item understand ramification, splitting, and inertness of rational primes;
    \item factor small rational primes in quadratic number fields;
    \item understand units and roots of unity in number fields;
    \item interpret Dirichlet's Unit Theorem conceptually;
    \item understand ideal class groups and class numbers;
    \item explain the relationship between principal ideals and unique factorization;
    \item recognize the role of discriminants;
    \item understand embeddings of number fields into \(\R\) and \(\C\);
    \item recognize how algebraic number theory connects to Diophantine equations,
    Galois theory, analytic number theory, and arithmetic geometry.
\end{itemize}
}

%%%%%%%%%%%%%%%%%%%%%%%%%%%%%%%%%%%%%%%%%%%%%%%%%%%%%%%%%%%%
\subsection{What Is Algebraic Number Theory?}
\label{subsec:what-is-algebraic-number-theory}
%%%%%%%%%%%%%%%%%%%%%%%%%%%%%%%%%%%%%%%%%%%%%%%%%%%%%%%%%%%%

Elementary number theory studies arithmetic primarily inside the ring of
integers

\[ \Z. \]

Algebraic number theory enlarges this setting.

Instead of studying only ordinary integers and rational numbers, we study
numbers obtained as solutions of polynomial equations with rational
coefficients.

\begin{definitionbox}{Algebraic Number Theory}
\emph{Algebraic number theory} is the study of algebraic numbers, number
fields, algebraic integers, and the arithmetic structures that arise inside
finite extensions of the rational numbers.
\end{definitionbox}

Its central philosophy is:

\[
\boxed{
\text{integer arithmetic}
\longrightarrow
\text{polynomial extensions}
\longrightarrow
\text{new arithmetic structures}.
}
\]

Many classical questions about ordinary integers become easier to understand
after moving to a larger number system.

%%%%%%%%%%%%%%%%%%%%%%%%%%%%%%%%%%%%%%%%%%%%%%%%%%%%%%%%%%%%
\subsection{Algebraic Numbers}
\label{subsec:algebraic-numbers}
%%%%%%%%%%%%%%%%%%%%%%%%%%%%%%%%%%%%%%%%%%%%%%%%%%%%%%%%%%%%

\begin{definitionbox}{Algebraic Number}
A complex number \(\alpha\in\C\) is called \emph{algebraic over \(\Q\)} if
there exists a nonzero polynomial

\[ f(x)\in\Q[x] \]

such that

\[ f(\alpha)=0. \]
\end{definitionbox}

The Greek letter

\[ \alpha \]

is called \emph{alpha}, pronounced ``AL-fuh.''

For example,

\[ \sqrt2 \]

is algebraic because it satisfies

\[ x^2-2=0. \]

Similarly,

\[ i \]

is algebraic because

\[ i^2+1=0. \]

Every rational number is algebraic.

Indeed, if

\[ r=\frac ab\in\Q, \]

then \(r\) satisfies

\[ bx-a=0. \]

%%%%%%%%%%%%%%%%%%%%%%%%%%%%%%%%%%%%%%%%%%%%%%%%%%%%%%%%%%%%
\subsection{Transcendental Numbers}
\label{subsec:transcendental-numbers}
%%%%%%%%%%%%%%%%%%%%%%%%%%%%%%%%%%%%%%%%%%%%%%%%%%%%%%%%%%%%

\begin{definitionbox}{Transcendental Number}
A complex number that is not algebraic over \(\Q\) is called
\emph{transcendental}.
\end{definitionbox}

Famous examples include

\[ \pi \]

and

\[ e. \]

Thus every complex number belongs to one of two broad classes:

\[
\boxed{
\C
=
\{\text{algebraic numbers}\}
\cup
\{\text{transcendental numbers}\}.
}
\]

%%%%%%%%%%%%%%%%%%%%%%%%%%%%%%%%%%%%%%%%%%%%%%%%%%%%%%%%%%%%
\subsection{The Field of Algebraic Numbers}
\label{subsec:field-algebraic-numbers}
%%%%%%%%%%%%%%%%%%%%%%%%%%%%%%%%%%%%%%%%%%%%%%%%%%%%%%%%%%%%

The algebraic numbers form a field.

That is, if \(\alpha\) and \(\beta\) are algebraic, then so are

\[ \alpha+\beta,\qquad\alpha-\beta,\qquad\alpha\beta, \]

and, if \(\beta\neq0\),

\[ \frac{\alpha}{\beta}. \]

This field is often denoted

\[ \overline{\Q}. \]

The notation

\[ \overline{\Q} \]

is read ``Q bar.''

It denotes the algebraic closure of \(\Q\) inside \(\C\).

%%%%%%%%%%%%%%%%%%%%%%%%%%%%%%%%%%%%%%%%%%%%%%%%%%%%%%%%%%%%
\subsection{Minimal Polynomials}
\label{subsec:minimal-polynomials-number-theory}
%%%%%%%%%%%%%%%%%%%%%%%%%%%%%%%%%%%%%%%%%%%%%%%%%%%%%%%%%%%%

An algebraic number may satisfy many polynomial equations.

For example, \(\sqrt2\) satisfies

\[ x^2-2=0 \]

but also

\[ x^4-4=0. \]

The most important polynomial is the minimal one.

\begin{definitionbox}{Minimal Polynomial}
Let \(\alpha\) be algebraic over \(\Q\). The \emph{minimal polynomial} of
\(\alpha\) over \(\Q\) is the unique monic irreducible polynomial

\[ m_\alpha(x)\in\Q[x] \]

such that

\[ m_\alpha(\alpha)=0. \]
\end{definitionbox}

The notation

\[ m_\alpha(x) \]

is read ``the minimal polynomial of alpha.''

For example,

\[ m_{\sqrt2}(x)=x^2-2. \]

%%%%%%%%%%%%%%%%%%%%%%%%%%%%%%%%%%%%%%%%%%%%%%%%%%%%%%%%%%%%
\subsection{Degree of an Algebraic Number}
\label{subsec:degree-algebraic-number}
%%%%%%%%%%%%%%%%%%%%%%%%%%%%%%%%%%%%%%%%%%%%%%%%%%%%%%%%%%%%

\begin{definitionbox}{Degree of an Algebraic Number}
The \emph{degree} of an algebraic number \(\alpha\) over \(\Q\) is the degree
of its minimal polynomial:

\[ [\Q(\alpha):\Q]=\deg m_\alpha(x). \]
\end{definitionbox}

For example,

\[ [\Q(\sqrt2):\Q]=2. \]

Likewise, if

\[ \alpha=\sqrt[3]{2}, \]

then

\[ m_\alpha(x)=x^3-2 \]

and

\[ [\Q(\sqrt[3]{2}):\Q]=3. \]

%%%%%%%%%%%%%%%%%%%%%%%%%%%%%%%%%%%%%%%%%%%%%%%%%%%%%%%%%%%%
\subsection{Number Fields}
\label{subsec:number-fields}
%%%%%%%%%%%%%%%%%%%%%%%%%%%%%%%%%%%%%%%%%%%%%%%%%%%%%%%%%%%%

\begin{definitionbox}{Number Field}
A \emph{number field} is a finite field extension

\[ K/\Q. \]

Equivalently, \(K\) is a field containing \(\Q\) such that

\[ [K:\Q]<\infty. \]
\end{definitionbox}

The symbol

\[ K/\Q \]

is read ``\(K\) over \(\Q\).''

Typical examples include

\[ \Q(\sqrt2),\qquad\Q(i),\qquad\Q(\sqrt[3]{2}). \]

%%%%%%%%%%%%%%%%%%%%%%%%%%%%%%%%%%%%%%%%%%%%%%%%%%%%%%%%%%%%
\subsection{Visualizing a Number Field Extension}
\label{subsec:number-field-extension-visual}
%%%%%%%%%%%%%%%%%%%%%%%%%%%%%%%%%%%%%%%%%%%%%%%%%%%%%%%%%%%%

\begin{centereddiagram}
\begin{tikzcd}[row sep=large]
K=\Q(\alpha) \arrow[d, dash] \\
\Q
\end{tikzcd}
\end{centereddiagram}

The extension degree

\[ [K:\Q]=n \]

means that \(K\) is an \(n\)-dimensional vector space over \(\Q\).

%%%%%%%%%%%%%%%%%%%%%%%%%%%%%%%%%%%%%%%%%%%%%%%%%%%%%%%%%%%%
\subsection{Worked Example: The Field \(\Q(\sqrt2)\)}
\label{subsec:worked-q-sqrt2}
%%%%%%%%%%%%%%%%%%%%%%%%%%%%%%%%%%%%%%%%%%%%%%%%%%%%%%%%%%%%

Every element of

\[ \Q(\sqrt2) \]

has the form

\[ a+b\sqrt2,\qquad a,b\in\Q. \]

Thus a basis over \(\Q\) is

\[ \{1,\sqrt2\}. \]

Therefore

\[ [\Q(\sqrt2):\Q]=2. \]

\begin{workedexamplebox}{Multiplication in a Quadratic Number Field}

Multiply

\[ (3+2\sqrt2)(1-\sqrt2). \]

\begin{tabularx}{\textwidth}{@{}>{\raggedright\arraybackslash}p{0.42\textwidth}X@{}}
Expand. &
\(3-3\sqrt2+2\sqrt2-4\)\\
Combine rational terms. &
\(3-4=-1\)\\
Combine irrational terms. &
\(-3\sqrt2+2\sqrt2=-\sqrt2\)
\end{tabularx}

\begin{answerbox}
\[ \boxed{(3+2\sqrt2)(1-\sqrt2)=-1-\sqrt2.} \]
\end{answerbox}
\end{workedexamplebox}

%%%%%%%%%%%%%%%%%%%%%%%%%%%%%%%%%%%%%%%%%%%%%%%%%%%%%%%%%%%%
\subsection{Quadratic Number Fields}
\label{subsec:quadratic-number-fields}
%%%%%%%%%%%%%%%%%%%%%%%%%%%%%%%%%%%%%%%%%%%%%%%%%%%%%%%%%%%%

A particularly important class of number fields has degree two.

\begin{definitionbox}{Quadratic Number Field}
A \emph{quadratic number field} is a field of the form

\[ \Q(\sqrt d), \]

where \(d\in\Z\) is squarefree and \(d\neq0,1\).
\end{definitionbox}

If

\[ d>0, \]

the field is called a \emph{real quadratic field}.

If

\[ d<0, \]

the field is called an \emph{imaginary quadratic field}.

Examples include

\[ \Q(\sqrt5),\qquad\Q(\sqrt{-1})=\Q(i),\qquad\Q(\sqrt{-5}). \]

%%%%%%%%%%%%%%%%%%%%%%%%%%%%%%%%%%%%%%%%%%%%%%%%%%%%%%%%%%%%
\subsection{Algebraic Integers}
\label{subsec:algebraic-integers}
%%%%%%%%%%%%%%%%%%%%%%%%%%%%%%%%%%%%%%%%%%%%%%%%%%%%%%%%%%%%

The ordinary integers are characterized as roots of monic polynomials with
integer coefficients.

This motivates the correct generalization of the integers inside a number
field.

\begin{definitionbox}{Algebraic Integer}
An algebraic number \(\alpha\) is an \emph{algebraic integer} if it satisfies
a monic polynomial

\[ f(x)\in\Z[x]. \]
\end{definitionbox}

For example,

\[ \sqrt2 \]

is an algebraic integer because it satisfies

\[ x^2-2=0. \]

Similarly,

\[ \frac{1+\sqrt5}{2} \]

is an algebraic integer because it satisfies

\[ x^2-x-1=0. \]

%%%%%%%%%%%%%%%%%%%%%%%%%%%%%%%%%%%%%%%%%%%%%%%%%%%%%%%%%%%%
\subsection{Not Every Algebraic Number Is an Algebraic Integer}
\label{subsec:algebraic-not-integer}
%%%%%%%%%%%%%%%%%%%%%%%%%%%%%%%%%%%%%%%%%%%%%%%%%%%%%%%%%%%%

Consider

\[ \frac12. \]

It is algebraic because it satisfies

\[ 2x-1=0. \]

However, the polynomial is not monic.

In fact, \(\frac12\) is not an algebraic integer.

Thus

\[
\boxed{
\text{algebraic integer}
\implies
\text{algebraic number},
}
\]

but the converse is false.

%%%%%%%%%%%%%%%%%%%%%%%%%%%%%%%%%%%%%%%%%%%%%%%%%%%%%%%%%%%%
\subsection{The Ring of Integers}
\label{subsec:ring-of-integers-number-field}
%%%%%%%%%%%%%%%%%%%%%%%%%%%%%%%%%%%%%%%%%%%%%%%%%%%%%%%%%%%%

\begin{definitionbox}{Ring of Integers}
Let \(K\) be a number field. The \emph{ring of integers} of \(K\), denoted

\[ \mathcal O_K, \]

is the set of all algebraic integers contained in \(K\).
\end{definitionbox}

The notation

\[ \mathcal O_K \]

is read ``O sub \(K\)'' or ``the ring of integers of \(K\).''

For the rational field,

\[ \mathcal O_{\Q}=\Z. \]

%%%%%%%%%%%%%%%%%%%%%%%%%%%%%%%%%%%%%%%%%%%%%%%%%%%%%%%%%%%%
\subsection{Quadratic Rings of Integers}
\label{subsec:quadratic-rings-of-integers}
%%%%%%%%%%%%%%%%%%%%%%%%%%%%%%%%%%%%%%%%%%%%%%%%%%%%%%%%%%%%

For a squarefree integer \(d\), the ring of integers of

\[ K=\Q(\sqrt d) \]

depends on \(d\pmod4\).

\begin{theorem}[Quadratic Ring of Integers]
Let \(d\) be squarefree. Then

\[
\mathcal O_K=
\begin{cases}
\Z\left[\dfrac{1+\sqrt d}{2}\right],
& d\equiv1\pmod4,\\[8pt]
\Z[\sqrt d],
& d\equiv2,3\pmod4.
\end{cases}
\]
\end{theorem}

For example,

\[ \mathcal O_{\Q(i)}=\Z[i]. \]

But for

\[ K=\Q(\sqrt5), \]

we have

\[ \mathcal O_K=\Z\left[\frac{1+\sqrt5}{2}\right]. \]

%%%%%%%%%%%%%%%%%%%%%%%%%%%%%%%%%%%%%%%%%%%%%%%%%%%%%%%%%%%%
\subsection{The Gaussian Integers}
\label{subsec:gaussian-integers}
%%%%%%%%%%%%%%%%%%%%%%%%%%%%%%%%%%%%%%%%%%%%%%%%%%%%%%%%%%%%

One of the most important examples is the ring

\[ \Z[i]=\{a+bi:a,b\in\Z\}. \]

Its elements are called \emph{Gaussian integers}.

\begin{centereddiagram}
\begin{tikzpicture}[scale=0.7]
\draw[-{Latex[length=2.5mm]}] (-4.5,0)--(4.5,0) node[right] {\(\Re\)};
\draw[-{Latex[length=2.5mm]}] (0,-4.5)--(0,4.5) node[above] {\(\Im\)};

\foreach \x in {-4,...,4}{
    \foreach \y in {-4,...,4}{
        \fill (\x,\y) circle (1.4pt);
    }
}

\node[above right] at (2,3) {\(2+3i\)};
\draw[fill] (2,3) circle (2.2pt);
\end{tikzpicture}
\end{centereddiagram}

The Gaussian integers form a square lattice in the complex plane.

%%%%%%%%%%%%%%%%%%%%%%%%%%%%%%%%%%%%%%%%%%%%%%%%%%%%%%%%%%%%
\subsection{Conjugation in Quadratic Fields}
\label{subsec:quadratic-conjugation}
%%%%%%%%%%%%%%%%%%%%%%%%%%%%%%%%%%%%%%%%%%%%%%%%%%%%%%%%%%%%

For

\[ \alpha=a+b\sqrt d, \]

define its conjugate by

\[ \overline{\alpha}=a-b\sqrt d. \]

For example,

\[ \overline{3+2\sqrt5}=3-2\sqrt5. \]

In \(\Q(i)\),

\[ \overline{a+bi}=a-bi, \]

which agrees with ordinary complex conjugation.

%%%%%%%%%%%%%%%%%%%%%%%%%%%%%%%%%%%%%%%%%%%%%%%%%%%%%%%%%%%%
\subsection{Field Trace}
\label{subsec:field-trace-number-theory}
%%%%%%%%%%%%%%%%%%%%%%%%%%%%%%%%%%%%%%%%%%%%%%%%%%%%%%%%%%%%

\begin{definitionbox}{Field Trace}
For a quadratic number field and

\[ \alpha=a+b\sqrt d, \]

the \emph{trace} is

\[ \operatorname{Tr}_{K/\Q}(\alpha)=\alpha+\overline{\alpha}. \]
\end{definitionbox}

Thus

\[ \operatorname{Tr}_{K/\Q}(a+b\sqrt d)=2a. \]

The notation

\[ \operatorname{Tr}_{K/\Q}(\alpha) \]

is read ``the trace from \(K\) to \(\Q\) of alpha.''

%%%%%%%%%%%%%%%%%%%%%%%%%%%%%%%%%%%%%%%%%%%%%%%%%%%%%%%%%%%%
\subsection{Field Norm}
\label{subsec:field-norm-number-theory}
%%%%%%%%%%%%%%%%%%%%%%%%%%%%%%%%%%%%%%%%%%%%%%%%%%%%%%%%%%%%

\begin{definitionbox}{Field Norm}
For a quadratic number field and

\[ \alpha=a+b\sqrt d, \]

the \emph{norm} is

\[ N_{K/\Q}(\alpha)=\alpha\overline{\alpha}. \]
\end{definitionbox}

Therefore

\[ \boxed{N_{K/\Q}(a+b\sqrt d)=a^2-db^2.} \]

For Gaussian integers,

\[ N(a+bi)=a^2+b^2. \]

%%%%%%%%%%%%%%%%%%%%%%%%%%%%%%%%%%%%%%%%%%%%%%%%%%%%%%%%%%%%
\subsection{Worked Example: Trace and Norm}
\label{subsec:worked-trace-norm}
%%%%%%%%%%%%%%%%%%%%%%%%%%%%%%%%%%%%%%%%%%%%%%%%%%%%%%%%%%%%

\begin{workedexamplebox}{Trace and Norm in \(\Q(\sqrt5)\)}

Let

\[ \alpha=3+2\sqrt5. \]

Find its trace and norm.

\begin{tabularx}{\textwidth}{@{}>{\raggedright\arraybackslash}p{0.42\textwidth}X@{}}
Find the conjugate. &
\(\overline{\alpha}=3-2\sqrt5\)\\
Add the conjugates. &
\(\operatorname{Tr}(\alpha)=6\)\\
Multiply the conjugates. &
\(N(\alpha)=(3+2\sqrt5)(3-2\sqrt5)\)\\
Simplify. &
\(N(\alpha)=9-20=-11\)
\end{tabularx}

\begin{answerbox}
\[ \boxed{\operatorname{Tr}(\alpha)=6,\qquad N(\alpha)=-11.} \]
\end{answerbox}
\end{workedexamplebox}

%%%%%%%%%%%%%%%%%%%%%%%%%%%%%%%%%%%%%%%%%%%%%%%%%%%%%%%%%%%%
\subsection{Multiplicativity of the Norm}
\label{subsec:norm-multiplicative}
%%%%%%%%%%%%%%%%%%%%%%%%%%%%%%%%%%%%%%%%%%%%%%%%%%%%%%%%%%%%

The norm satisfies

\[ \boxed{N(\alpha\beta)=N(\alpha)N(\beta).} \]

This makes the norm extremely useful in factorization problems.

For Gaussian integers,

\[ N(a+bi)=a^2+b^2\geq0. \]

If

\[ \alpha\beta=1, \]

then

\[ N(\alpha)N(\beta)=1. \]

Therefore units must have norm \(1\).

%%%%%%%%%%%%%%%%%%%%%%%%%%%%%%%%%%%%%%%%%%%%%%%%%%%%%%%%%%%%
\subsection{Units}
\label{subsec:number-field-units}
%%%%%%%%%%%%%%%%%%%%%%%%%%%%%%%%%%%%%%%%%%%%%%%%%%%%%%%%%%%%

\begin{definitionbox}{Unit}
An element \(u\in\mathcal O_K\) is a \emph{unit} if there exists
\(v\in\mathcal O_K\) such that

\[ uv=1. \]
\end{definitionbox}

For Gaussian integers,

\[ N(a+bi)=1 \]

implies

\[ a^2+b^2=1. \]

Thus the Gaussian units are

\[ \boxed{\pm1,\pm i.} \]

%%%%%%%%%%%%%%%%%%%%%%%%%%%%%%%%%%%%%%%%%%%%%%%%%%%%%%%%%%%%
\subsection{Associates}
\label{subsec:associates-number-fields}
%%%%%%%%%%%%%%%%%%%%%%%%%%%%%%%%%%%%%%%%%%%%%%%%%%%%%%%%%%%%

Two elements \(\alpha,\beta\in\mathcal O_K\) are called
\emph{associates} if

\[ \alpha=u\beta \]

for some unit \(u\).

For example, in \(\Z[i]\),

\[ 1+i,\quad -1-i,\quad -1+i,\quad1-i \]

are all associates.

Factorization is therefore considered unique only up to multiplication by
units and reordering of factors.

%%%%%%%%%%%%%%%%%%%%%%%%%%%%%%%%%%%%%%%%%%%%%%%%%%%%%%%%%%%%
\subsection{Irreducible and Prime Elements}
\label{subsec:irreducible-prime-elements}
%%%%%%%%%%%%%%%%%%%%%%%%%%%%%%%%%%%%%%%%%%%%%%%%%%%%%%%%%%%%

In \(\Z\), prime and irreducible elements coincide.

In more general rings, they need not.

\begin{definitionbox}{Irreducible Element}
A nonzero nonunit \(\alpha\) is \emph{irreducible} if

\[ \alpha=\beta\gamma \]

implies that either \(\beta\) or \(\gamma\) is a unit.
\end{definitionbox}

\begin{definitionbox}{Prime Element}
A nonzero nonunit \(\pi\) is \emph{prime} if

\[ \pi\mid\alpha\beta \]

implies

\[ \pi\mid\alpha\qquad\text{or}\qquad\pi\mid\beta. \]
\end{definitionbox}

Every prime element is irreducible in an integral domain, but the converse
can fail.

%%%%%%%%%%%%%%%%%%%%%%%%%%%%%%%%%%%%%%%%%%%%%%%%%%%%%%%%%%%%
\subsection{Failure of Unique Factorization}
\label{subsec:failure-unique-factorization}
%%%%%%%%%%%%%%%%%%%%%%%%%%%%%%%%%%%%%%%%%%%%%%%%%%%%%%%%%%%%

The ring

\[ \Z[\sqrt{-5}] \]

provides a classical example.

Consider the integer \(6\).

We have

\[ 6=2\cdot3. \]

But also,

\[ 6=(1+\sqrt{-5})(1-\sqrt{-5}). \]

These factorizations are genuinely different.

In this ring, the factors involved can be shown to be irreducible, but they
are not associates.

Thus unique factorization into irreducible elements fails.

\begin{warningbox}
The Fundamental Theorem of Arithmetic is special to rings with sufficient
factorization properties.

When the ambient ring changes, unique factorization of elements may fail.
\end{warningbox}

%%%%%%%%%%%%%%%%%%%%%%%%%%%%%%%%%%%%%%%%%%%%%%%%%%%%%%%%%%%%
\subsection{Why Ideals Solve the Problem}
\label{subsec:why-ideals-solve-factorization}
%%%%%%%%%%%%%%%%%%%%%%%%%%%%%%%%%%%%%%%%%%%%%%%%%%%%%%%%%%%%

The central achievement of algebraic number theory is that even when
elements do not factor uniquely, ideals do.

Instead of factoring

\[ \alpha \]

into elements, we factor the principal ideal

\[ (\alpha). \]

This restores a form of unique factorization.

%%%%%%%%%%%%%%%%%%%%%%%%%%%%%%%%%%%%%%%%%%%%%%%%%%%%%%%%%%%%
\subsection{Ideals in Number Rings}
\label{subsec:ideals-number-rings}
%%%%%%%%%%%%%%%%%%%%%%%%%%%%%%%%%%%%%%%%%%%%%%%%%%%%%%%%%%%%

The general concept of an ideal is developed in Ring Theory.

For a ring of integers \(\mathcal O_K\), an ideal

\[ I\subseteq\mathcal O_K \]

is an additive subgroup satisfying

\[ rI\subseteq I \]

for every

\[ r\in\mathcal O_K. \]

A principal ideal generated by \(\alpha\) is

\[ (\alpha)=\{\alpha r:r\in\mathcal O_K\}. \]

%%%%%%%%%%%%%%%%%%%%%%%%%%%%%%%%%%%%%%%%%%%%%%%%%%%%%%%%%%%%
\subsection{Prime Ideals}
\label{subsec:prime-ideals-number-theory}
%%%%%%%%%%%%%%%%%%%%%%%%%%%%%%%%%%%%%%%%%%%%%%%%%%%%%%%%%%%%

\begin{definitionbox}{Prime Ideal}
A proper ideal \(\mathfrak p\subsetneq\mathcal O_K\) is called
\emph{prime} if

\[ ab\in\mathfrak p \]

implies

\[ a\in\mathfrak p\qquad\text{or}\qquad b\in\mathfrak p. \]
\end{definitionbox}

The symbol

\[ \mathfrak p \]

is a lowercase Fraktur \(p\), commonly used to denote prime ideals.

%%%%%%%%%%%%%%%%%%%%%%%%%%%%%%%%%%%%%%%%%%%%%%%%%%%%%%%%%%%%
\subsection{Unique Factorization of Ideals}
\label{subsec:unique-factorization-ideals}
%%%%%%%%%%%%%%%%%%%%%%%%%%%%%%%%%%%%%%%%%%%%%%%%%%%%%%%%%%%%

\begin{theorem}[Unique Factorization of Ideals]
Let \(K\) be a number field and let \(\mathcal O_K\) be its ring of integers.

Every nonzero proper ideal \(I\subseteq\mathcal O_K\) can be written uniquely
in the form

\[ \boxed{I=\mathfrak p_1^{e_1}\cdots\mathfrak p_r^{e_r}}, \]

where the \(\mathfrak p_i\) are distinct nonzero prime ideals and

\[ e_i\geq1. \]
\end{theorem}

This is the ideal-theoretic replacement for the Fundamental Theorem of
Arithmetic.

%%%%%%%%%%%%%%%%%%%%%%%%%%%%%%%%%%%%%%%%%%%%%%%%%%%%%%%%%%%%
\subsection{Dedekind Domains}
\label{subsec:dedekind-domains-number-theory}
%%%%%%%%%%%%%%%%%%%%%%%%%%%%%%%%%%%%%%%%%%%%%%%%%%%%%%%%%%%%

Rings in which ideals behave this well belong to an important class.

\begin{definitionbox}{Dedekind Domain}
A \emph{Dedekind domain} is an integral domain in which every nonzero proper
ideal factors uniquely into nonzero prime ideals, together with equivalent
structural conditions developed in Commutative Algebra.
\end{definitionbox}

The ring of integers

\[ \mathcal O_K \]

of every number field \(K\) is a Dedekind domain.

%%%%%%%%%%%%%%%%%%%%%%%%%%%%%%%%%%%%%%%%%%%%%%%%%%%%%%%%%%%%
\subsection{Fractional Ideals}
\label{subsec:fractional-ideals}
%%%%%%%%%%%%%%%%%%%%%%%%%%%%%%%%%%%%%%%%%%%%%%%%%%%%%%%%%%%%

Ordinary ideals can be generalized by allowing denominators.

\begin{definitionbox}{Fractional Ideal}
A \emph{fractional ideal} of a number field \(K\) is an
\(\mathcal O_K\)-submodule \(I\subseteq K\) such that there exists a nonzero

\[ d\in\mathcal O_K \]

with

\[ dI\subseteq\mathcal O_K. \]
\end{definitionbox}

Fractional ideals form a multiplicative group under ideal multiplication.

This allows ideal quotients and inverses to be treated naturally.

%%%%%%%%%%%%%%%%%%%%%%%%%%%%%%%%%%%%%%%%%%%%%%%%%%%%%%%%%%%%
\subsection{Ideal Classes}
\label{subsec:ideal-classes}
%%%%%%%%%%%%%%%%%%%%%%%%%%%%%%%%%%%%%%%%%%%%%%%%%%%%%%%%%%%%

Two nonzero fractional ideals \(I\) and \(J\) are considered equivalent if

\[ I=\alpha J \]

for some nonzero

\[ \alpha\in K. \]

The equivalence classes form the ideal class group.

\begin{definitionbox}{Ideal Class Group}
The \emph{ideal class group} of a number field \(K\), denoted

\[ \operatorname{Cl}(K), \]

is the quotient of the group of nonzero fractional ideals by the subgroup
of principal fractional ideals.
\end{definitionbox}

The notation

\[ \operatorname{Cl}(K) \]

is read ``the class group of \(K\).''

%%%%%%%%%%%%%%%%%%%%%%%%%%%%%%%%%%%%%%%%%%%%%%%%%%%%%%%%%%%%
\subsection{Class Number}
\label{subsec:class-number}
%%%%%%%%%%%%%%%%%%%%%%%%%%%%%%%%%%%%%%%%%%%%%%%%%%%%%%%%%%%%

\begin{definitionbox}{Class Number}
The \emph{class number} of \(K\) is

\[ h_K=\#\operatorname{Cl}(K). \]
\end{definitionbox}

The class number measures the failure of every ideal to be principal.

If

\[ h_K=1, \]

then every ideal is principal.

In rings of integers of number fields, this implies that the ring is a
principal ideal domain and therefore a unique factorization domain.

%%%%%%%%%%%%%%%%%%%%%%%%%%%%%%%%%%%%%%%%%%%%%%%%%%%%%%%%%%%%
\subsection{Factorization of Rational Primes}
\label{subsec:rational-prime-factorization}
%%%%%%%%%%%%%%%%%%%%%%%%%%%%%%%%%%%%%%%%%%%%%%%%%%%%%%%%%%%%

A rational prime \(p\in\Z\) may behave differently when regarded inside a
larger ring of integers.

The ideal

\[ (p)\subseteq\mathcal O_K \]

may:

\begin{itemize}
    \item remain prime;
    \item split into distinct prime ideals;
    \item ramify, producing repeated prime-ideal factors.
\end{itemize}

These possibilities are fundamental to algebraic number theory.

%%%%%%%%%%%%%%%%%%%%%%%%%%%%%%%%%%%%%%%%%%%%%%%%%%%%%%%%%%%%
\subsection{Splitting, Inertness, and Ramification}
\label{subsec:splitting-inert-ramification}
%%%%%%%%%%%%%%%%%%%%%%%%%%%%%%%%%%%%%%%%%%%%%%%%%%%%%%%%%%%%

In a quadratic number field, a rational prime \(p\) can typically have one of
three behaviors.

\begin{definitionbox}{Split Prime}
A rational prime \(p\) \emph{splits} if

\[ (p)=\mathfrak p_1\mathfrak p_2 \]

with distinct prime ideals

\[ \mathfrak p_1\neq\mathfrak p_2. \]
\end{definitionbox}

\begin{definitionbox}{Inert Prime}
A rational prime \(p\) is \emph{inert} if the ideal

\[ (p) \]

remains prime in \(\mathcal O_K\).
\end{definitionbox}

\begin{definitionbox}{Ramified Prime}
A rational prime \(p\) is \emph{ramified} if

\[ (p)=\mathfrak p^2 \]

in a quadratic number field.
\end{definitionbox}

%%%%%%%%%%%%%%%%%%%%%%%%%%%%%%%%%%%%%%%%%%%%%%%%%%%%%%%%%%%%
\subsection{Prime Behavior Diagram}
\label{subsec:prime-behavior-diagram}
%%%%%%%%%%%%%%%%%%%%%%%%%%%%%%%%%%%%%%%%%%%%%%%%%%%%%%%%%%%%

\begin{centereddiagram}
\begin{tikzpicture}[
    node distance=12mm,
    box/.style={draw,rounded corners,align=center,minimum width=3.3cm,
    minimum height=0.9cm}
]
\node[box] (p) {Rational prime \(p\)};
\node[box,below left=of p] (split) {Split\\\((p)=\mathfrak p_1\mathfrak p_2\)};
\node[box,below=of p] (inert) {Inert\\\((p)\) remains prime};
\node[box,below right=of p] (ram) {Ramified\\\((p)=\mathfrak p^2\)};

\draw[-{Latex[length=2.5mm]}] (p)--(split);
\draw[-{Latex[length=2.5mm]}] (p)--(inert);
\draw[-{Latex[length=2.5mm]}] (p)--(ram);
\end{tikzpicture}
\end{centereddiagram}

%%%%%%%%%%%%%%%%%%%%%%%%%%%%%%%%%%%%%%%%%%%%%%%%%%%%%%%%%%%%
\subsection{Prime Splitting in the Gaussian Integers}
\label{subsec:prime-splitting-gaussian}
%%%%%%%%%%%%%%%%%%%%%%%%%%%%%%%%%%%%%%%%%%%%%%%%%%%%%%%%%%%%

In

\[ \Z[i], \]

rational primes exhibit a particularly elegant pattern.

\begin{theorem}[Prime Behavior in \(\Z[i]\)]
Let \(p\) be an odd rational prime.

\begin{itemize}
    \item If
    \[ p\equiv1\pmod4, \]
    then \(p\) splits in \(\Z[i]\).

    \item If
    \[ p\equiv3\pmod4, \]
    then \(p\) remains prime in \(\Z[i]\).

    \item The prime \(2\) ramifies:
    \[ 2=-i(1+i)^2. \]
\end{itemize}
\end{theorem}

%%%%%%%%%%%%%%%%%%%%%%%%%%%%%%%%%%%%%%%%%%%%%%%%%%%%%%%%%%%%
\subsection{Worked Example: Splitting of \(5\)}
\label{subsec:worked-splitting-five}
%%%%%%%%%%%%%%%%%%%%%%%%%%%%%%%%%%%%%%%%%%%%%%%%%%%%%%%%%%%%

Since

\[ 5\equiv1\pmod4, \]

the prime \(5\) splits in \(\Z[i]\).

Indeed,

\[ 5=(2+i)(2-i). \]

Because

\[ N(2+i)=2^2+1^2=5, \]

the factors have prime norm.

\begin{answerbox}
\[ \boxed{5=(2+i)(2-i)\text{ in }\Z[i].} \]
\end{answerbox}

%%%%%%%%%%%%%%%%%%%%%%%%%%%%%%%%%%%%%%%%%%%%%%%%%%%%%%%%%%%%
\subsection{Worked Example: The Prime \(3\)}
\label{subsec:worked-inert-three}
%%%%%%%%%%%%%%%%%%%%%%%%%%%%%%%%%%%%%%%%%%%%%%%%%%%%%%%%%%%%

Since

\[ 3\equiv3\pmod4, \]

the prime \(3\) remains prime in \(\Z[i]\).

There are no integers \(a,b\) satisfying

\[ a^2+b^2=3. \]

Thus \(3\) cannot factor into nonunits of smaller positive norm.

%%%%%%%%%%%%%%%%%%%%%%%%%%%%%%%%%%%%%%%%%%%%%%%%%%%%%%%%%%%%
\subsection{Connection with Sums of Two Squares}
\label{subsec:sums-two-squares-connection}
%%%%%%%%%%%%%%%%%%%%%%%%%%%%%%%%%%%%%%%%%%%%%%%%%%%%%%%%%%%%

The splitting of primes in \(\Z[i]\) is closely related to representation as
a sum of two squares.

\begin{theorem}[Fermat's Two-Square Theorem]
An odd prime \(p\) can be written as

\[ p=a^2+b^2 \]

with integers \(a,b\) if and only if

\[ p\equiv1\pmod4. \]
\end{theorem}

For example,

\[ 5=1^2+2^2, \]

\[ 13=2^2+3^2, \]

\[ 17=1^2+4^2. \]

This theorem demonstrates how factorization in an enlarged ring can solve an
integer representation problem.

%%%%%%%%%%%%%%%%%%%%%%%%%%%%%%%%%%%%%%%%%%%%%%%%%%%%%%%%%%%%
\subsection{Integral Bases}
\label{subsec:integral-bases}
%%%%%%%%%%%%%%%%%%%%%%%%%%%%%%%%%%%%%%%%%%%%%%%%%%%%%%%%%%%%

The ring of integers \(\mathcal O_K\) is a free \(\Z\)-module of rank

\[ n=[K:\Q]. \]

Thus there exist algebraic integers

\[ \omega_1,\ldots,\omega_n \]

such that every

\[ \alpha\in\mathcal O_K \]

can be written uniquely as

\[ \alpha=a_1\omega_1+\cdots+a_n\omega_n, \]

where

\[ a_i\in\Z. \]

\begin{definitionbox}{Integral Basis}
A basis

\[ \{\omega_1,\ldots,\omega_n\} \]

of \(\mathcal O_K\) as a free \(\Z\)-module is called an
\emph{integral basis}.
\end{definitionbox}

%%%%%%%%%%%%%%%%%%%%%%%%%%%%%%%%%%%%%%%%%%%%%%%%%%%%%%%%%%%%
\subsection{Embeddings of Number Fields}
\label{subsec:number-field-embeddings}
%%%%%%%%%%%%%%%%%%%%%%%%%%%%%%%%%%%%%%%%%%%%%%%%%%%%%%%%%%%%

A number field can be mapped into \(\C\) in several algebraically compatible
ways.

\begin{definitionbox}{Embedding}
An \emph{embedding} of a number field \(K\) into \(\C\) is an injective field
homomorphism

\[ \sigma:K\to\C \]

that fixes every element of \(\Q\).
\end{definitionbox}

The Greek letter

\[ \sigma \]

is called \emph{sigma}.

For

\[ K=\Q(\sqrt2), \]

there are two embeddings:

\[ \sigma_1(\sqrt2)=\sqrt2 \]

and

\[ \sigma_2(\sqrt2)=-\sqrt2. \]

%%%%%%%%%%%%%%%%%%%%%%%%%%%%%%%%%%%%%%%%%%%%%%%%%%%%%%%%%%%%
\subsection{Real and Complex Embeddings}
\label{subsec:real-complex-embeddings}
%%%%%%%%%%%%%%%%%%%%%%%%%%%%%%%%%%%%%%%%%%%%%%%%%%%%%%%%%%%%

An embedding

\[ \sigma:K\to\C \]

is called \emph{real} if

\[ \sigma(K)\subseteq\R. \]

Otherwise, embeddings occur in complex-conjugate pairs.

If \(K\) has

\[ r_1 \]

real embeddings and

\[ r_2 \]

pairs of complex embeddings, then

\[ \boxed{[K:\Q]=r_1+2r_2.} \]

The pair

\[ (r_1,r_2) \]

is called the \emph{signature} of the number field.

%%%%%%%%%%%%%%%%%%%%%%%%%%%%%%%%%%%%%%%%%%%%%%%%%%%%%%%%%%%%
\subsection{Trace and Norm via Embeddings}
\label{subsec:trace-norm-embeddings}
%%%%%%%%%%%%%%%%%%%%%%%%%%%%%%%%%%%%%%%%%%%%%%%%%%%%%%%%%%%%

For a degree-\(n\) number field with embeddings

\[ \sigma_1,\ldots,\sigma_n, \]

the trace and norm may be written as

\[ \boxed{\operatorname{Tr}_{K/\Q}(\alpha)=\sum_{i=1}^{n}\sigma_i(\alpha)} \]

and

\[ \boxed{N_{K/\Q}(\alpha)=\prod_{i=1}^{n}\sigma_i(\alpha).} \]

This generalizes the quadratic formulas introduced earlier.

%%%%%%%%%%%%%%%%%%%%%%%%%%%%%%%%%%%%%%%%%%%%%%%%%%%%%%%%%%%%
\subsection{Discriminants}
\label{subsec:number-field-discriminants}
%%%%%%%%%%%%%%%%%%%%%%%%%%%%%%%%%%%%%%%%%%%%%%%%%%%%%%%%%%%%

The discriminant measures important arithmetic information about a number
field.

Let

\[ \omega_1,\ldots,\omega_n \]

be a basis of \(K/\Q\).

\begin{definitionbox}{Discriminant of a Basis}
The discriminant of the basis is

\[
\operatorname{disc}(\omega_1,\ldots,\omega_n)
=
\det\left(
\operatorname{Tr}_{K/\Q}(\omega_i\omega_j)
\right)_{i,j=1}^{n}.
\]
\end{definitionbox}

The determinant notation

\[ \det(\cdot) \]

was introduced in Linear Algebra.

For an integral basis, this value is the discriminant of the number field.

%%%%%%%%%%%%%%%%%%%%%%%%%%%%%%%%%%%%%%%%%%%%%%%%%%%%%%%%%%%%
\subsection{Quadratic Field Discriminants}
\label{subsec:quadratic-field-discriminants}
%%%%%%%%%%%%%%%%%%%%%%%%%%%%%%%%%%%%%%%%%%%%%%%%%%%%%%%%%%%%

For a squarefree integer \(d\), the discriminant of

\[ K=\Q(\sqrt d) \]

is

\[
D_K=
\begin{cases}
d, & d\equiv1\pmod4,\\
4d, & d\equiv2,3\pmod4.
\end{cases}
\]

For example,

\[ D_{\Q(i)}=-4 \]

and

\[ D_{\Q(\sqrt5)}=5. \]

%%%%%%%%%%%%%%%%%%%%%%%%%%%%%%%%%%%%%%%%%%%%%%%%%%%%%%%%%%%%
\subsection{Ramification and the Discriminant}
\label{subsec:ramification-discriminant}
%%%%%%%%%%%%%%%%%%%%%%%%%%%%%%%%%%%%%%%%%%%%%%%%%%%%%%%%%%%%

The discriminant detects ramification.

\begin{theorem}
If a rational prime \(p\) ramifies in a number field \(K\), then

\[ p\mid D_K. \]
\end{theorem}

For quadratic fields, the converse also holds.

Thus in

\[ \Q(i), \]

the discriminant is

\[ -4, \]

so the only ramified rational prime is

\[ 2. \]

%%%%%%%%%%%%%%%%%%%%%%%%%%%%%%%%%%%%%%%%%%%%%%%%%%%%%%%%%%%%
\subsection{Norm of an Ideal}
\label{subsec:ideal-norm}
%%%%%%%%%%%%%%%%%%%%%%%%%%%%%%%%%%%%%%%%%%%%%%%%%%%%%%%%%%%%

The norm concept extends from elements to ideals.

\begin{definitionbox}{Ideal Norm}
For a nonzero ideal

\[ I\subseteq\mathcal O_K, \]

its norm is

\[ N(I)=\#(\mathcal O_K/I). \]
\end{definitionbox}

Thus \(N(I)\) is the number of residue classes modulo \(I\).

Ideal norms satisfy

\[ \boxed{N(IJ)=N(I)N(J).} \]

For a principal ideal generated by \(\alpha\),

\[ N((\alpha))=|N_{K/\Q}(\alpha)|. \]

%%%%%%%%%%%%%%%%%%%%%%%%%%%%%%%%%%%%%%%%%%%%%%%%%%%%%%%%%%%%
\subsection{Prime Ideal Decomposition}
\label{subsec:prime-ideal-decomposition}
%%%%%%%%%%%%%%%%%%%%%%%%%%%%%%%%%%%%%%%%%%%%%%%%%%%%%%%%%%%%

For a rational prime \(p\), its ideal factorization has the form

\[ (p)=\mathfrak p_1^{e_1}\cdots\mathfrak p_g^{e_g}. \]

Each prime ideal also has a residue degree

\[ f_i=[\mathcal O_K/\mathfrak p_i:\mathbf F_p]. \]

The notation

\[ \mathbf F_p \]

denotes the finite field with \(p\) elements.

For number fields,

\[ \boxed{\sum_{i=1}^{g}e_i f_i=[K:\Q]} \]

under the standard prime-factorization framework.

The integers \(e_i\) are called \emph{ramification indices}, and the
integers \(f_i\) are called \emph{residue degrees}.

%%%%%%%%%%%%%%%%%%%%%%%%%%%%%%%%%%%%%%%%%%%%%%%%%%%%%%%%%%%%
\subsection{Units in General Number Fields}
\label{subsec:units-general-number-fields}
%%%%%%%%%%%%%%%%%%%%%%%%%%%%%%%%%%%%%%%%%%%%%%%%%%%%%%%%%%%%

The unit group

\[ \mathcal O_K^\times \]

can be much larger than the units of \(\Z\).

The notation

\[ \mathcal O_K^\times \]

is read ``the unit group of \(\mathcal O_K\).''

In \(\Z\),

\[ \Z^\times=\{\pm1\}. \]

In \(\Z[i]\),

\[ \Z[i]^\times=\{\pm1,\pm i\}. \]

In real quadratic fields, there are infinitely many units.

%%%%%%%%%%%%%%%%%%%%%%%%%%%%%%%%%%%%%%%%%%%%%%%%%%%%%%%%%%%%
\subsection{Dirichlet's Unit Theorem}
\label{subsec:dirichlet-unit-theorem}
%%%%%%%%%%%%%%%%%%%%%%%%%%%%%%%%%%%%%%%%%%%%%%%%%%%%%%%%%%%%

\begin{theorem}[Dirichlet's Unit Theorem]
Let \(K\) be a number field with \(r_1\) real embeddings and \(r_2\) pairs of
complex embeddings.

Then

\[ \boxed{\mathcal O_K^\times\cong\mu_K\times\Z^{r_1+r_2-1}}, \]

where \(\mu_K\) is the finite group of roots of unity contained in \(K\).
\end{theorem}

Here

\[ \mu_K \]

is read ``mu sub \(K\).''

This theorem gives the structure of the unit group.

The number

\[ r_1+r_2-1 \]

is called the \emph{unit rank}.

%%%%%%%%%%%%%%%%%%%%%%%%%%%%%%%%%%%%%%%%%%%%%%%%%%%%%%%%%%%%
\subsection{Example: Units in \(\Q(\sqrt2)\)}
\label{subsec:units-q-sqrt2}
%%%%%%%%%%%%%%%%%%%%%%%%%%%%%%%%%%%%%%%%%%%%%%%%%%%%%%%%%%%%

The field

\[ \Q(\sqrt2) \]

has

\[ r_1=2,\qquad r_2=0. \]

Therefore its unit rank is

\[ 2+0-1=1. \]

A fundamental unit is

\[ 1+\sqrt2. \]

Its norm is

\[ N(1+\sqrt2)=1-2=-1. \]

Powers of this element generate infinitely many units up to sign.

%%%%%%%%%%%%%%%%%%%%%%%%%%%%%%%%%%%%%%%%%%%%%%%%%%%%%%%%%%%%
\subsection{Pell Equations and Units}
\label{subsec:pell-equations-units}
%%%%%%%%%%%%%%%%%%%%%%%%%%%%%%%%%%%%%%%%%%%%%%%%%%%%%%%%%%%%

Consider the equation

\[ x^2-dy^2=1. \]

This is a Pell equation.

If

\[ \alpha=x+y\sqrt d, \]

then

\[ N(\alpha)=x^2-dy^2. \]

Therefore solutions of

\[ x^2-dy^2=1 \]

correspond to units of norm \(1\) in

\[ \Z[\sqrt d] \]

or the appropriate ring of integers.

\begin{conceptbox}
A Diophantine equation can become a statement about units:

\[
\boxed{
x^2-dy^2=1
\quad\Longleftrightarrow\quad
N(x+y\sqrt d)=1.
}
\]

This is a typical algebraic-number-theory strategy:
translate an integer equation into arithmetic inside a number field.
\end{conceptbox}

%%%%%%%%%%%%%%%%%%%%%%%%%%%%%%%%%%%%%%%%%%%%%%%%%%%%%%%%%%%%
\subsection{Worked Example: Pell Equation}
\label{subsec:worked-pell-equation}
%%%%%%%%%%%%%%%%%%%%%%%%%%%%%%%%%%%%%%%%%%%%%%%%%%%%%%%%%%%%

\begin{workedexamplebox}{A Solution of a Pell Equation}

Verify that

\[ x=3,\qquad y=2 \]

solves

\[ x^2-2y^2=1. \]

\begin{tabularx}{\textwidth}{@{}>{\raggedright\arraybackslash}p{0.42\textwidth}X@{}}
Substitute. &
\(3^2-2(2^2)\)\\
Evaluate. &
\(9-8=1\)
\end{tabularx}

Equivalently,

\[ N(3+2\sqrt2)=1. \]

\begin{answerbox}
\[ \boxed{(x,y)=(3,2)\text{ is a solution.}} \]
\end{answerbox}
\end{workedexamplebox}

%%%%%%%%%%%%%%%%%%%%%%%%%%%%%%%%%%%%%%%%%%%%%%%%%%%%%%%%%%%%
\subsection{Minkowski's Geometric Viewpoint}
\label{subsec:minkowski-viewpoint}
%%%%%%%%%%%%%%%%%%%%%%%%%%%%%%%%%%%%%%%%%%%%%%%%%%%%%%%%%%%%

A number field may be embedded into a real vector space using all of its
embeddings simultaneously.

For a field with signature

\[ (r_1,r_2), \]

one obtains a map into

\[ \R^{r_1}\times\C^{r_2}. \]

Since

\[ \C\cong\R^2, \]

this may be viewed as a lattice in

\[ \R^n, \]

where

\[ n=[K:\Q]. \]

This geometric viewpoint allows methods from lattice theory and convex
geometry to be applied to algebraic number theory.

%%%%%%%%%%%%%%%%%%%%%%%%%%%%%%%%%%%%%%%%%%%%%%%%%%%%%%%%%%%%
\subsection{Minkowski Embedding}
\label{subsec:minkowski-embedding}
%%%%%%%%%%%%%%%%%%%%%%%%%%%%%%%%%%%%%%%%%%%%%%%%%%%%%%%%%%%%

Let the embeddings of \(K\) be organized as real embeddings and one from
each complex-conjugate pair.

The map

\[
\Phi:K\to\R^{r_1}\times\C^{r_2}
\]

defined by

\[
\Phi(\alpha)
=
\left(
\sigma_1(\alpha),\ldots,\sigma_{r_1+r_2}(\alpha)
\right)
\]

is called a \emph{Minkowski embedding}.

Under an appropriate real-coordinate identification,

\[ \Phi(\mathcal O_K) \]

forms a lattice.

%%%%%%%%%%%%%%%%%%%%%%%%%%%%%%%%%%%%%%%%%%%%%%%%%%%%%%%%%%%%
\subsection{Why Class Groups Are Finite}
\label{subsec:class-group-finiteness}
%%%%%%%%%%%%%%%%%%%%%%%%%%%%%%%%%%%%%%%%%%%%%%%%%%%%%%%%%%%%

A fundamental theorem states that

\[ \operatorname{Cl}(K) \]

is finite for every number field \(K\).

The proof uses geometry of numbers and Minkowski's theorem.

Thus the class number

\[ h_K \]

is always a finite positive integer.

This finiteness is one of the major structural results of algebraic number
theory.

%%%%%%%%%%%%%%%%%%%%%%%%%%%%%%%%%%%%%%%%%%%%%%%%%%%%%%%%%%%%
\subsection{Class Number One}
\label{subsec:class-number-one}
%%%%%%%%%%%%%%%%%%%%%%%%%%%%%%%%%%%%%%%%%%%%%%%%%%%%%%%%%%%%

When

\[ h_K=1, \]

every ideal is principal.

Consequently, the ring of integers is a principal ideal domain and hence a
unique factorization domain.

For example,

\[ \Z[i] \]

has class number \(1\).

By contrast,

\[ \Z[\sqrt{-5}] \]

has nontrivial class group, explaining why unique factorization of elements
fails.

%%%%%%%%%%%%%%%%%%%%%%%%%%%%%%%%%%%%%%%%%%%%%%%%%%%%%%%%%%%%
\subsection{Cyclotomic Fields}
\label{subsec:cyclotomic-fields-introduction}
%%%%%%%%%%%%%%%%%%%%%%%%%%%%%%%%%%%%%%%%%%%%%%%%%%%%%%%%%%%%

Roots of unity generate an important class of number fields.

Let

\[ \zeta_n=e^{2\pi i/n}. \]

The symbol

\[ \zeta_n \]

is read ``zeta sub \(n\).''

It is a primitive \(n\)-th root of unity.

\begin{definitionbox}{Cyclotomic Field}
The field

\[ \Q(\zeta_n) \]

is called the \emph{\(n\)-th cyclotomic field}.
\end{definitionbox}

Cyclotomic fields connect:

\begin{itemize}
    \item roots of unity;
    \item Galois theory;
    \item congruences;
    \item prime splitting;
    \item reciprocity laws.
\end{itemize}

%%%%%%%%%%%%%%%%%%%%%%%%%%%%%%%%%%%%%%%%%%%%%%%%%%%%%%%%%%%%
\subsection{Cyclotomic Polynomials}
\label{subsec:cyclotomic-polynomials-number-theory}
%%%%%%%%%%%%%%%%%%%%%%%%%%%%%%%%%%%%%%%%%%%%%%%%%%%%%%%%%%%%

The minimal polynomial of a primitive \(n\)-th root of unity is the
cyclotomic polynomial

\[ \Phi_n(x). \]

The symbol

\[ \Phi \]

is uppercase Greek \emph{phi}.

The polynomial satisfies

\[ x^n-1=\prod_{d\mid n}\Phi_d(x). \]

Its degree is

\[ \boxed{\deg\Phi_n=\varphi(n).} \]

Therefore

\[ [\Q(\zeta_n):\Q]=\varphi(n). \]

%%%%%%%%%%%%%%%%%%%%%%%%%%%%%%%%%%%%%%%%%%%%%%%%%%%%%%%%%%%%
\subsection{Example: Fourth Cyclotomic Field}
\label{subsec:fourth-cyclotomic-field}
%%%%%%%%%%%%%%%%%%%%%%%%%%%%%%%%%%%%%%%%%%%%%%%%%%%%%%%%%%%%

For

\[ n=4, \]

a primitive fourth root of unity is

\[ \zeta_4=i. \]

Its cyclotomic polynomial is

\[ \Phi_4(x)=x^2+1. \]

Therefore

\[ \Q(\zeta_4)=\Q(i). \]

This recovers the Gaussian number field as a cyclotomic field.

%%%%%%%%%%%%%%%%%%%%%%%%%%%%%%%%%%%%%%%%%%%%%%%%%%%%%%%%%%%%
\subsection{Prime Splitting and Polynomial Factorization}
\label{subsec:prime-splitting-polynomial}
%%%%%%%%%%%%%%%%%%%%%%%%%%%%%%%%%%%%%%%%%%%%%%%%%%%%%%%%%%%%

The splitting of rational primes in number fields is closely connected to
factorization of defining polynomials modulo \(p\).

Suppose

\[ K=\Q(\alpha) \]

and \(\alpha\) satisfies a polynomial

\[ f(x)\in\Z[x]. \]

Factoring \(f(x)\) modulo \(p\) often gives information about how

\[ (p) \]

factors in \(\mathcal O_K\).

This principle becomes precise through results such as Dedekind's
factorization theorem.

%%%%%%%%%%%%%%%%%%%%%%%%%%%%%%%%%%%%%%%%%%%%%%%%%%%%%%%%%%%%
\subsection{Algebraic Number Theory and Diophantine Equations}
\label{subsec:algebraic-number-theory-diophantine}
%%%%%%%%%%%%%%%%%%%%%%%%%%%%%%%%%%%%%%%%%%%%%%%%%%%%%%%%%%%%

A polynomial equation in integers can sometimes be factored more naturally
inside a number field.

For example,

\[ x^2+y^2=z^2 \]

is connected to factorization in the Gaussian integers because

\[ x^2+y^2=(x+iy)(x-iy). \]

Likewise,

\[ x^2-dy^2=1 \]

becomes a norm equation.

More generally,

\[
\boxed{
\text{Diophantine equation}
\longrightarrow
\text{factorization or norm equation in a number field}.
}
\]

%%%%%%%%%%%%%%%%%%%%%%%%%%%%%%%%%%%%%%%%%%%%%%%%%%%%%%%%%%%%
\subsection{Algebraic Number Theory and Galois Theory}
\label{subsec:algebraic-number-theory-galois}
%%%%%%%%%%%%%%%%%%%%%%%%%%%%%%%%%%%%%%%%%%%%%%%%%%%%%%%%%%%%

Galois theory studies field extensions through their automorphism groups.

For a Galois number field extension

\[ K/\Q, \]

the group

\[ \operatorname{Gal}(K/\Q) \]

acts on algebraic integers, ideals, and prime ideals.

Prime splitting is strongly influenced by this group action.

This connection eventually leads to Frobenius elements and one of the major
results of modern number theory: the Chebotarev Density Theorem.

%%%%%%%%%%%%%%%%%%%%%%%%%%%%%%%%%%%%%%%%%%%%%%%%%%%%%%%%%%%%
\subsection{Algebraic Number Theory and Analytic Number Theory}
\label{subsec:algebraic-analytic-number-theory}
%%%%%%%%%%%%%%%%%%%%%%%%%%%%%%%%%%%%%%%%%%%%%%%%%%%%%%%%%%%%

The Riemann zeta function

\[ \zeta(s) \]

encodes arithmetic information about \(\Z\).

A number field has its own zeta function.

\begin{definitionbox}{Dedekind Zeta Function}
For a number field \(K\), the \emph{Dedekind zeta function} is defined in its
initial region of convergence by

\[ \boxed{\zeta_K(s)=\sum_{I\neq0}\frac{1}{N(I)^s}}, \]

where the sum runs over nonzero ideals of \(\mathcal O_K\).
\end{definitionbox}

It has an Euler product

\[ \boxed{\zeta_K(s)=\prod_{\mathfrak p}\frac{1}{1-N(\mathfrak p)^{-s}}}. \]

Thus prime ideals replace ordinary prime numbers.

%%%%%%%%%%%%%%%%%%%%%%%%%%%%%%%%%%%%%%%%%%%%%%%%%%%%%%%%%%%%
\subsection{The Riemann Zeta Function as a Special Case}
\label{subsec:riemann-zeta-special-case}
%%%%%%%%%%%%%%%%%%%%%%%%%%%%%%%%%%%%%%%%%%%%%%%%%%%%%%%%%%%%

When

\[ K=\Q, \]

the ring of integers is

\[ \mathcal O_{\Q}=\Z. \]

Every nonzero ideal has the form

\[ (n). \]

Hence

\[ N((n))=n. \]

Therefore

\[ \zeta_{\Q}(s)=\zeta(s). \]

So the Riemann zeta function is the Dedekind zeta function of the rational
number field.

%%%%%%%%%%%%%%%%%%%%%%%%%%%%%%%%%%%%%%%%%%%%%%%%%%%%%%%%%%%%
\subsection{Common Errors}
\label{subsec:algebraic-number-theory-common-errors}
%%%%%%%%%%%%%%%%%%%%%%%%%%%%%%%%%%%%%%%%%%%%%%%%%%%%%%%%%%%%

\subsubsection{Assuming Every Irrational Number Is Transcendental}

Many irrational numbers are algebraic.

For example,

\[ \sqrt2 \]

is irrational but algebraic.

%%%%%%%%%%%%%%%%%%%%%%%%%%%%%%%%%%%%%%%%%%%%%%%%%%%%%%%%%%%%
\subsubsection{Confusing Algebraic Numbers with Algebraic Integers}

An algebraic integer must satisfy a monic polynomial with integer
coefficients.

An algebraic number only needs to satisfy a nonzero polynomial with rational
coefficients.

%%%%%%%%%%%%%%%%%%%%%%%%%%%%%%%%%%%%%%%%%%%%%%%%%%%%%%%%%%%%
\subsubsection{Assuming \(\mathcal O_{\Q(\sqrt d)}=\Z[\sqrt d]\) Always}

This fails when

\[ d\equiv1\pmod4. \]

In that case,

\[ \mathcal O_K=\Z\left[\frac{1+\sqrt d}{2}\right]. \]

%%%%%%%%%%%%%%%%%%%%%%%%%%%%%%%%%%%%%%%%%%%%%%%%%%%%%%%%%%%%
\subsubsection{Assuming Unique Factorization of Elements Always Holds}

Unique factorization can fail in rings of integers.

What always survives is unique factorization of nonzero ideals into prime
ideals.

%%%%%%%%%%%%%%%%%%%%%%%%%%%%%%%%%%%%%%%%%%%%%%%%%%%%%%%%%%%%
\subsubsection{Confusing Irreducible and Prime Elements}

In a general integral domain, irreducible does not necessarily imply prime.

%%%%%%%%%%%%%%%%%%%%%%%%%%%%%%%%%%%%%%%%%%%%%%%%%%%%%%%%%%%%
\subsubsection{Confusing Element Norm and Ideal Norm}

For an element,

\[ N_{K/\Q}(\alpha) \]

is defined through embeddings or conjugates.

For an ideal,

\[ N(I)=\#(\mathcal O_K/I). \]

The concepts are related but distinct.

%%%%%%%%%%%%%%%%%%%%%%%%%%%%%%%%%%%%%%%%%%%%%%%%%%%%%%%%%%%%
\subsubsection{Assuming a Rational Prime Stays Prime}

A rational prime may split, remain inert, or ramify after passing to a larger
number field.

%%%%%%%%%%%%%%%%%%%%%%%%%%%%%%%%%%%%%%%%%%%%%%%%%%%%%%%%%%%%
\subsubsection{Confusing \(\Z[\sqrt d]\) with \(\Q(\sqrt d)\)}

The ring

\[ \Z[\sqrt d] \]

contains integer combinations of \(1\) and \(\sqrt d\).

The field

\[ \Q(\sqrt d) \]

contains rational combinations and is closed under division by nonzero
elements.

%%%%%%%%%%%%%%%%%%%%%%%%%%%%%%%%%%%%%%%%%%%%%%%%%%%%%%%%%%%%
\subsubsection{Assuming Class Number One Means Only One Ideal}

The statement

\[ h_K=1 \]

means there is only one ideal class.

There are still infinitely many ideals.

%%%%%%%%%%%%%%%%%%%%%%%%%%%%%%%%%%%%%%%%%%%%%%%%%%%%%%%%%%%%
\subsubsection{Ignoring Units in Factorization}

Factorizations are considered equivalent up to multiplication by units and
reordering.

%%%%%%%%%%%%%%%%%%%%%%%%%%%%%%%%%%%%%%%%%%%%%%%%%%%%%%%%%%%%
\subsection{Applications and Connections}
\label{subsec:algebraic-number-theory-applications}
%%%%%%%%%%%%%%%%%%%%%%%%%%%%%%%%%%%%%%%%%%%%%%%%%%%%%%%%%%%%

\begin{applicationbox}
\textbf{Diophantine Equations.}
Algebraic extensions can convert polynomial equations into factorization,
ideal, or norm problems.

\medskip

\textbf{Galois Theory.}
Automorphisms of number fields control symmetries among algebraic numbers and
influence the behavior of prime ideals.

\medskip

\textbf{Commutative Algebra.}
Rings of integers are Dedekind domains, and their ideals provide a central
example of localization, integral closure, prime ideals, and factorization.

\medskip

\textbf{Analytic Number Theory.}
Dedekind zeta functions and more general \(L\)-functions encode prime-ideal
distribution analytically.

\medskip

\textbf{Geometry of Numbers.}
Minkowski embeddings turn rings of integers and ideals into lattices,
allowing geometric methods to prove arithmetic results.

\medskip

\textbf{Cryptography.}
Arithmetic in number fields, elliptic curves, and ideal class groups appears
in modern cryptographic constructions.

\medskip

\textbf{Arithmetic Geometry.}
Number fields provide the base arithmetic setting for the study of algebraic
curves, elliptic curves, and higher-dimensional varieties.

\medskip

\textbf{Representation Problems.}
Questions such as representing primes as sums of squares are naturally
explained by factorization in suitable number rings.
\end{applicationbox}

%%%%%%%%%%%%%%%%%%%%%%%%%%%%%%%%%%%%%%%%%%%%%%%%%%%%%%%%%%%%
\subsection{Exercises}
\label{subsec:algebraic-number-theory-exercises}
%%%%%%%%%%%%%%%%%%%%%%%%%%%%%%%%%%%%%%%%%%%%%%%%%%%%%%%%%%%%

\begin{exercise}[1]
Show that

\[ \sqrt3 \]

is algebraic over \(\Q\).
\end{exercise}

\begin{exercise}[1]
Find a polynomial in \(\Z[x]\) satisfied by

\[ \sqrt[3]{5}. \]
\end{exercise}

\begin{exercise}[2]
Find the minimal polynomial over \(\Q\) of

\[ 1+\sqrt2. \]
\end{exercise}

\begin{exercise}[2]
Find the minimal polynomial over \(\Q\) of

\[ \sqrt2+\sqrt3. \]
\end{exercise}

\begin{exercise}[1]
Determine

\[ [\Q(\sqrt7):\Q]. \]
\end{exercise}

\begin{exercise}[2]
Write the general element of

\[ \Q(\sqrt{-3}). \]
\end{exercise}

\begin{exercise}[2]
Show that

\[ \frac{1+\sqrt5}{2} \]

is an algebraic integer.
\end{exercise}

\begin{exercise}[2]
Explain why

\[ \frac13 \]

is algebraic but not an algebraic integer.
\end{exercise}

\begin{exercise}[1]
Determine the ring of integers of

\[ \Q(\sqrt2). \]
\end{exercise}

\begin{exercise}[1]
Determine the ring of integers of

\[ \Q(\sqrt{13}). \]
\end{exercise}

\begin{exercise}[2]
For

\[ \alpha=4+3\sqrt2, \]

compute the conjugate, trace, and norm.
\end{exercise}

\begin{exercise}[2]
For

\[ \alpha=3+4i, \]

compute

\[ N(\alpha). \]
\end{exercise}

\begin{exercise}[2]
Verify directly that

\[ N(\alpha\beta)=N(\alpha)N(\beta) \]

for two general Gaussian integers.
\end{exercise}

\begin{exercise}[1]
List all units in

\[ \Z[i]. \]
\end{exercise}

\begin{exercise}[2]
Determine whether

\[ 2+i \]

is a unit in \(\Z[i]\).
\end{exercise}

\begin{exercise}[2]
Factor

\[ 13 \]

in \(\Z[i]\).
\end{exercise}

\begin{exercise}[2]
Explain why \(7\) remains prime in \(\Z[i]\).
\end{exercise}

\begin{exercise}[2]
Show that

\[ 17=1^2+4^2 \]

and use this to factor \(17\) in \(\Z[i]\).
\end{exercise}

\begin{exercise}[3]
Verify the two factorizations

\[ 6=2\cdot3 \]

and

\[ 6=(1+\sqrt{-5})(1-\sqrt{-5}) \]

inside \(\Z[\sqrt{-5}]\).
\end{exercise}

\begin{exercise}[3]
Explain why the existence of two different factorizations in
\(\Z[\sqrt{-5}]\) indicates a failure of unique factorization of elements.
\end{exercise}

\begin{exercise}[2]
State the unique factorization theorem for ideals in the ring of integers of
a number field.
\end{exercise}

\begin{exercise}[2]
Explain the difference among a split, inert, and ramified rational prime.
\end{exercise}

\begin{exercise}[1]
Determine whether each odd prime splits or remains inert in \(\Z[i]\):

\[ 3,\qquad5,\qquad7,\qquad13,\qquad19. \]
\end{exercise}

\begin{exercise}[2]
Find the discriminant of

\[ \Q(\sqrt2). \]
\end{exercise}

\begin{exercise}[2]
Find the discriminant of

\[ \Q(\sqrt5). \]
\end{exercise}

\begin{exercise}[2]
Use the discriminant to determine which rational primes can ramify in

\[ \Q(\sqrt5). \]
\end{exercise}

\begin{exercise}[2]
Find the two embeddings

\[ \Q(\sqrt3)\to\R. \]
\end{exercise}

\begin{exercise}[2]
Determine the signature

\[ (r_1,r_2) \]

of

\[ \Q(\sqrt2). \]
\end{exercise}

\begin{exercise}[2]
Determine the signature of

\[ \Q(i). \]
\end{exercise}

\begin{exercise}[2]
Use Dirichlet's Unit Theorem to determine the unit rank of

\[ \Q(\sqrt5). \]
\end{exercise}

\begin{exercise}[2]
Use Dirichlet's Unit Theorem to determine the unit rank of

\[ \Q(i). \]
\end{exercise}

\begin{exercise}[2]
Verify that

\[ 1+\sqrt2 \]

is a unit in \(\Z[\sqrt2]\).
\end{exercise}

\begin{exercise}[2]
Verify that

\[ 3+2\sqrt2 \]

has norm \(1\).
\end{exercise}

\begin{exercise}[2]
Use the norm interpretation to verify that

\[ (x,y)=(3,2) \]

solves the Pell equation

\[ x^2-2y^2=1. \]
\end{exercise}

\begin{exercise}[2]
Compute

\[ \Phi_3(x),\qquad\Phi_4(x),\qquad\Phi_6(x). \]
\end{exercise}

\begin{exercise}[2]
Determine

\[ [\Q(\zeta_5):\Q]. \]
\end{exercise}

\begin{exercise}[2]
Explain why

\[ \Q(i) \]

is a cyclotomic field.
\end{exercise}

\begin{exercise}[3]
Explain how prime splitting in \(\Z[i]\) is related to the theorem on primes
represented as sums of two squares.
\end{exercise}

\begin{exercise}[3]
Explain conceptually why unique factorization of ideals can survive even when
unique factorization of elements fails.
\end{exercise}

\begin{exercise}[3]
Explain what the class number measures.
\end{exercise}

\begin{exercise}[3]
Explain why

\[ h_K=1 \]

implies that every ideal of \(\mathcal O_K\) is principal.
\end{exercise}

\begin{exercise}[3]
Explain the conceptual chain

\[
\text{polynomial equation}
\longrightarrow
\text{algebraic number}
\longrightarrow
\text{number field}
\longrightarrow
\text{ring of integers}
\longrightarrow
\text{ideal factorization}.
\]
\end{exercise}

\begin{exercise}[3]
Explain the conceptual relationship

\[
\text{rational primes}
\longrightarrow
\text{prime ideals}
\longrightarrow
\text{splitting behavior}
\longrightarrow
\text{arithmetic of the number field}.
\]
\end{exercise}

%%%%%%%%%%%%%%%%%%%%%%%%%%%%%%%%%%%%%%%%%%%%%%%%%%%%%%%%%%%%
\subsection{Conceptual Summary}
\label{subsec:algebraic-number-theory-summary}
%%%%%%%%%%%%%%%%%%%%%%%%%%%%%%%%%%%%%%%%%%%%%%%%%%%%%%%%%%%%

Algebraic number theory extends ordinary integer arithmetic into finite
extensions of the rational numbers.

An algebraic number is a complex number satisfying a polynomial equation over
\(\Q\).

A number field is a finite extension

\[ K/\Q. \]

The algebraic integers in \(K\) form its ring of integers

\[ \mathcal O_K. \]

Quadratic fields

\[ \Q(\sqrt d) \]

provide the simplest nontrivial examples.

Trace and norm generalize important arithmetic operations and connect
elements with their conjugates and embeddings.

Unique factorization of elements may fail in rings of integers, as in

\[ \Z[\sqrt{-5}]. \]

However, nonzero ideals factor uniquely into prime ideals:

\[ I=\mathfrak p_1^{e_1}\cdots\mathfrak p_r^{e_r}. \]

The ideal class group

\[ \operatorname{Cl}(K) \]

measures the failure of ideals to be principal, while the class number

\[ h_K \]

measures the size of this obstruction.

Rational primes can split, remain inert, or ramify in number fields.

The discriminant helps detect ramification.

Embeddings place number fields into real and complex spaces and lead to the
Minkowski geometric viewpoint.

Dirichlet's Unit Theorem describes the structure of

\[ \mathcal O_K^\times. \]

Cyclotomic fields connect roots of unity with Galois theory and prime
splitting.

Dedekind zeta functions extend the Riemann zeta function from ordinary
integers to ideals in number fields.

Together these ideas reveal a central theme:

\[
\boxed{
\text{arithmetic becomes more structured when integers are studied inside
larger algebraic systems}.
}
\]

%%%%%%%%%%%%%%%%%%%%%%%%%%%%%%%%%%%%%%%%%%%%%%%%%%%%%%%%%%%%
\subsection{Section Connections}
\label{subsec:algebraic-number-theory-connections}
%%%%%%%%%%%%%%%%%%%%%%%%%%%%%%%%%%%%%%%%%%%%%%%%%%%%%%%%%%%%

\chapterconnection{
Algebraic Number Theory builds directly on Elementary Number Theory and on
the algebraic structures developed in Chapter 1. Algebraic numbers lead to
number fields, while algebraic integers form rings whose ideals restore a
form of unique factorization even when element factorization fails.
Quadratic fields provide the first examples of prime splitting, ramification,
norms, units, and class groups. Galois Theory explains symmetries of field
extensions and later controls prime splitting through Frobenius elements.
Commutative Algebra provides the abstract framework for integral closure,
Dedekind domains, localization, and ideal theory. Analytic Number Theory
extends naturally through Dedekind zeta functions and other \(L\)-functions.
Norm equations and unit groups connect directly with Diophantine equations,
while class groups, elliptic curves, and higher-dimensional varieties lead
toward Diophantine Geometry and Arithmetic Geometry.
}

%%%%%%%%%%%%%%%%%%%%%%%%%%%%%%%%%%%%%%%%%%%%%%%%%%%%%%%%%%%%
% SECTION SOURCE TRACKING
%%%%%%%%%%%%%%%%%%%%%%%%%%%%%%%%%%%%%%%%%%%%%%%%%%%%%%%%%%%%

%------------------------------------------------------------
% 2.3 Algebraic Number Theory
%------------------------------------------------------------
%
% Source basis:
%   - Standard undergraduate and introductory graduate
%     algebraic number theory.
%
% Used for:
%   - algebraic and transcendental numbers;
%   - minimal polynomials and degrees;
%   - number fields;
%   - quadratic number fields;
%   - algebraic integers and rings of integers;
%   - Gaussian integers;
%   - conjugation, trace, and norm;
%   - units and associates;
%   - irreducible and prime elements;
%   - failure of unique factorization;
%   - ideal factorization;
%   - Dedekind domains;
%   - fractional ideals;
%   - class groups and class numbers;
%   - prime splitting, inertness, and ramification;
%   - discriminants;
%   - ideal norms;
%   - embeddings and signatures;
%   - Dirichlet's Unit Theorem;
%   - Pell equations;
%   - Minkowski embeddings;
%   - cyclotomic fields and cyclotomic polynomials;
%   - prime splitting via polynomial factorization;
%   - Dedekind zeta functions;
%   - connections to Diophantine equations, Galois Theory,
%     Analytic Number Theory, and Arithmetic Geometry.
%
% Notes:
%   - General field-extension theory is taught canonically in
%     Field Theory.
%   - Galois groups and fundamental Galois correspondence are
%     taught canonically in Galois Theory.
%   - General ideal theory, localization, integral dependence,
%     and Dedekind-domain structure are developed in
%     Commutative Algebra.
%   - Detailed zeta and L-function analysis belongs in
%     Analytic Number Theory.
%   - Advanced Diophantine methods are deferred to
%     Diophantine Geometry.
%   - Higher-dimensional arithmetic constructions are deferred
%     to Arithmetic Geometry.
%
%%%%%%%%%%%%%%%%%%%%%%%%%%%%%%%%%%%%%%%%%%%%%%%%%%%%%%%%%%%%